\section{Conclusion}
\label{sec.conclusion}

The $k$-deep rule is a theorem about validators. A validator is synchronised
with the other validators, so the chain it holds is never far behind the
longest honest chain, and a constant $k$ suffices for it. A client has no such
guarantee. The chain it receives can be short, and how short is exactly how
much of it the adversary could have built alone.

This paper gives the rule a client should use instead. It prunes
$\frac{t}{n-t}$ of the gap between the height of the chain it received and the
tallest height any chain can have reached by the current round. The output is a
prefix of every chain held by an honest validator from that round onwards, the
outputs of any two clients are prefixes of one another, and the output grows at
rate $\left(1-\frac{t^2}{(n-t)^2}\right)f$. None of this needs a bound on the
delay of the client's own link, and a client that never hears an honest
message is still safe. Private mining attains the coefficient, and no constant
depth is sound.

The price is a confirmation lag of $\frac{t^2}{(n-t)^2}$ of the age of the
chain, paid even by a client that receives an honest chain in the round it was
adopted. That is the cost of not being able to tell that the peer is honest.
The main question we leave open is the same rule under variable difficulty,
where height is no longer a proxy for work.
